\begin{lemma}
\label{lemma-flat-ML-criterion}
Let $M$ be a flat $R$-module. The following are equivalent
\begin{enumerate}
\item $M$ is Mittag-Leffler, and
\item if $F$ is a finite free $R$-module and
$x \in F \otimes_R M$, then there exists a smallest submodule $F'$ of $F$
such that $x \in F' \otimes_R M$.
\end{enumerate}
\end{lemma}

\begin{proof}
The implication (1) $\Rightarrow$ (2) is a special case of
Lemma \ref{lemma-minimal-contains}. Assume (2).
By Theorem \ref{theorem-lazard} we can write $M$ as the colimit
$M = \colim_{i \in I} M_i$ of a directed system $(M_i, f_{ij})$ of
finite free $R$-modules.
By Remark \ref{remark-flat-ML}, it suffices to show that the inverse system
$(\Hom_R(M_i, R), \Hom_R(f_{ij}, R))$ is Mittag-Leffler.  In
other words,
fix $i \in I$ and for $j \geq i$ let $Q_j$ be the image of
$\Hom_R(M_j, R)
\to \Hom_R(M_i, R)$; we must show that the $Q_j$ stabilize.

\medskip\noindent
Since $M_i$ is free and finite, we can make the identification
$\Hom_R(M_i, M_j) = \Hom_R(M_i, R) \otimes_R  M_j$ for all $j$.
 Using the
fact that the $M_j$ are free, it follows that for $j \geq i$, $Q_j$ is the
smallest submodule of $\Hom_R(M_i, R)$ such that $f_{ij} \in Q_j
\otimes_R
M_j$.  Under the identification $\Hom_R(M_i, M) = \Hom_R(M_i, R)
\otimes_R
M$, the canonical map $f_i: M_i \to M$ is in $\Hom_R(M_i, R)
\otimes_R M$.  By the assumption on $M$, there exists a smallest submodule
$Q$ of $\Hom_R(M_i, R)$ such that $f_i \in Q \otimes_R M$.  We are
going to
show that the $Q_j$ stabilize to $Q$.

\medskip\noindent
For $j \geq i$ we have a commutative diagram
$$
\xymatrix{
Q_j \otimes_R M_j \ar[r] \ar[d] & \Hom_R(M_i, R) \otimes_R M_j
\ar[d] \\
Q_j \otimes_R M \ar[r] & \Hom_R(M_i, R) \otimes_R M.
}
$$
Since $f_{ij} \in Q_j \otimes_R M_j$ maps to $f_i \in \Hom_R(M_i, R)
\otimes_R M$, it follows that $f_i \in Q_j \otimes_R M$.  Hence, by the
choice of $Q$, we have $Q \subset Q_j$ for all $j \geq i$.

\medskip\noindent
Since the $Q_j$ are decreasing and $Q \subset Q_j$ for all $j \geq i$, to show
that the $Q_j$ stabilize to $Q$ it suffices to find a $j \geq i$ such that $Q_j
\subset Q$.  As an element of
$$
\Hom_R(M_i, R) \otimes_R M = \colim_{j \in J}
(\Hom_R(M_i, R) \otimes_R
M_j),
$$
$f_i$ is the colimit of $f_{ij}$ for $j \geq i$, and $f_i$ also lies in the
submodule
$$
\colim_{j \in J} (Q \otimes_R M_j) \subset \colim_{j \in J}
(\Hom_R(M_i, R)
\otimes_R M_j).
$$
It follows that for some $j \geq i$, $f_{ij}$ lies in $Q \otimes_R M_j$.
Since $Q_j$ is the smallest submodule of $\Hom_R(M_i, R)$ with $f_{ij}
\in
Q_j \otimes_R M_j$, we conclude $Q_j\subset Q$.
\end{proof}
